% order=104 分段重写草稿；尚未完成，不纳入 rewritten 计数。
% 已完成范围：章首语、sec:ontology_reconstruction、sec:swarm_dynamics_equation。
% 续写位置：sec:phase_transition_trigger 起。
\chapter{流体群脑工程：分布式状态场、相态调度与可验证协作}
\label{chp:fluid_swarm_engineering}

当机器人数量从若干台扩展到数百乃至数千台时，困难并不只是通信量变大，而是控制对象发生了变化：工程师既要追踪每个具身节点的安全边界、能力与故障，又要判断群体覆盖、连通、拥塞和任务供需等整体量如何随时间演化。我把“流体群脑”定义为处理这一双重尺度问题的计算模型，而不是把机器人宣称为真正的流体粒子。局部结构网络保留可审计的节点身份与邻接关系，全局分布表示提供可压缩的宏观读出；二者通过任务相关绑定、核估计和控制接口联合。本章将相态理解为可测工作区间，将所谓场力重建为受约束控制项，并以时延、丢包、模型失配和执行器极限说明其适用边界。HSF-HD在这里给出智能的一般状态动力学语言，AgentOS只是其中一种工程实现，不构成所有群体智能的唯一机制。

\section{双尺度本体：具身节点、结构网络与分布状态}
\label{sec:ontology_reconstruction}

先给出对象和空间。设群体含有有限节点集 $V_t=\{1,\ldots,N_t\}$，节点数允许因加入、退出或失效而变化。机器人 $i$ 的物理状态为 $x_i(t)=(r_i,v_i,q_i)\in\mathcal X_i$，其中位置 $r_i$ 以米计，速度 $v_i$ 以米每秒计，姿态 $q_i$ 为无量纲旋转表示；资源与能力状态 $c_i(t)\in\mathbb R^{d_c}$ 记录剩余电量、载荷、传感器可用性和计算预算，各维必须在进入距离计算前按量纲归一化；任务角色状态 $z_i(t)\in\Delta^{d_z-1}$ 是角色概率单纯形上的向量，而不是不可验证的“角色本质”。节点的完整局部状态写成 $s_i=(x_i,c_i,z_i,\Sigma_i,\tau_i)$，其中 $\Sigma_i$ 是估计不确定度，$\tau_i$ 是最近一次可靠观测时间戳。这个定义把事实、估计和控制意图分开：位置测量是观测，电量是设备状态，角色向量是调度变量，三者不能用一个含混的“能量”名称代替。

局部协作结构由带属性动态图 $G_t=(V_t,E_t,A_t)$ 表示。边 $(i,j)\in E_t$ 只说明在给定通信与感知门槛下两节点可发生作用，边属性 $A_{ij}$ 则保存时延、带宽、相对位姿置信度、信任上限与任务依赖。节点身份在图中是显式的，因此碰撞责任、消息来源和故障隔离都可以追溯。所谓“前锋”“中继”或“搬运者”不是固定拓扑位置，而是由任务 $g_t$、能力约束和当前证据共同决定的绑定 $b_t:i\mapsto z_i$。当目标变化时，绑定也应进入状态转移；不能因节点曾处于队形前端，就把它永久等同于某个功能模块。局部图尤其适合表达稀疏接触、互斥关系和必须逐节点满足的安全约束，但在大群体中直接对全图求解会受计算量和通信半径限制。

为得到宏观读出，在物理区域 $\Omega\subset\mathbb R^d$ 与角色空间 $\mathcal Z$ 上定义核化经验测度
\[
 \rho_t(r,z)=\sum_{i\in V_t}w_i(t)K_{h_r}(r-r_i(t))K_{h_z}(z-z_i(t)).
\]
这里 $K_{h_r}$ 和 $K_{h_z}$ 是带宽分别为 $h_r$ 与 $h_z$ 的归一化核，$w_i\ge0$ 是按传感可靠度或可用能力设定的权重。若 $w_i=1$，对 $r,z$ 积分得到节点数 $N_t$；若权重表示归一化能力，则积分得到的只是能力总量，不能再解释为概率密度。全局分布可读出覆盖缺口、角色供需差、拥塞峰值和连通风险的近似量，却会丢失节点身份和某些高阶相关。两个相距很远的小队与一个连续长队可能产生相近的低阶密度，因此不得从 $\rho_t$ 单独推出通信图连通，更不能由“密度平滑”推出任务已经成功。

局部结构与全局分布通过两个方向的接口联合。上行接口 $\mathcal A$ 从带时间戳的局部状态聚合为 $\hat\rho_t=\mathcal A(\{s_i\},G_t)$，同时输出采样覆盖和估计误差界；下行接口 $\mathcal D$ 把宏观缺口、风险预算与任务优先级转换为节点可执行的参考量 $u_i^{ref}=\mathcal D(i,G_t,\hat\rho_t,g_t)$。这不是把一个“场”神秘地施加给个体，而是明示身份、绑定、耦合和读出规则。若网络分区，上行估计退化为各连通分量的局部版本，下行指令只能在已有证据覆盖内有效；系统必须以版本号和租约阻止多个过期宏观状态同时支配同一节点。AgentOS实例可以把 $s_i$ 的短时工作状态映射到 $h_i(t)$，把任务偏差映射到 $E_i$，再由 $\Theta_i$ 表示局部更新门槛，但这些符号只是实现选择。HSF-HD要求的是状态、观测、约束和更新的可区分性，并不要求所有群体都采用同一组模块。

这一双尺度定义也限定了“群脑”一词的含义。它表示多个自治单元在共享任务和有限通信下形成可观测的联合动力学，不表示存在一个与单体大脑同构的实体。脑区类比不能证明无线网络具有神经机制，流体类比也不能证明密度方程足以描述离散碰撞。模型选择应由误差和任务决定：在通道狭窄、节点稀疏或个体异构显著时，以图模型为主；在大范围覆盖、统计调度和负载均衡时，以分布读出为主；在协同搬运等强耦合任务中，两者必须同时保留。反例是两台关键承重机器人同时掉线：宏观密度可能几乎不变，结构图上的承载割集却已断裂。只有把割集约束送入分布层的风险读出，宏观调度才不会被表面平滑度误导。

\section{受约束群体动力学：从场语言到可执行控制}
\label{sec:swarm_dynamics_equation}

在动力学层面，我不把目的论名称当作方程来源，而从可测状态、控制输入和扰动开始。对质量为 $m_i$ 千克的移动机器人，连续时间模型写为
\[
 \dot r_i=v_i,\qquad
 m_i\dot v_i=u_i^{safe}+u_i^{task}+u_i^{coh}+u_i^{exp}+d_i,
\]
其中各 $u$ 项单位均为牛顿，$d_i$ 汇集未建模摩擦、风和执行器误差。若底盘接受速度而不是力命令，则应改用相应的一阶或非完整运动学模型，不能为了保留“力”的比喻而伪造质量参数。假设状态估计误差有界、控制周期上界已知、执行器满足 $\|u_i\|\le u_i^{max}$，并且在紧急制动距离内至少有一种独立近场感知。所谓认知洛伦兹力在此只保留为历史命名；工程上它分解为安全投影、任务梯度、邻域一致性和受预算约束的探索扰动。

安全项优先于其他目标。令障碍和节点对的安全函数分别为 $h_{io}(r_i)$ 与 $h_{ij}(r_i,r_j)$，正值代表满足最小间距。控制器在每个周期求解
\[
 u_i^{safe}=\arg\min_{u\in\mathcal U_i}\|u-u_i^{nom}\|^2,
 \quad\text{s.t.}\quad \dot h_{ik}+\alpha(h_{ik})\ge-\epsilon_{ik},
\]
其中 $u_i^{nom}=u_i^{task}+u_i^{coh}+u_i^{exp}$，$\alpha$ 为选定的增益函数，$\epsilon_{ik}\ge0$ 是在约束冲突时显式付费的松弛量。这个二次规划是工程设计；在模型和采样条件满足时，它可以给出前向不变性的数学结果，但有感知漏检、制动饱和或求解超时时只能声称风险降低，不能声称绝对安全。原稿把排斥力类比为不相容原理，这种说法在解释层可以帮助直觉，却不提供任何安全证明；真正的保证来自安全集合、时延裕量和故障制动路径。

任务项由随时间变化的代价 $J_i(r,z;g_t,\hat\rho_t)$ 产生，典型选择为 $u_i^{task}=-k_g\nabla_{r_i}J_i$。如果目标 $g_t$、分布估计 $\hat\rho_t$ 或地图参数也在变化，沿轨迹的导数必须写成
\[
 \frac{dJ_i}{dt}=\nabla_{r_i}J_i^\top v_i+
 \partial_gJ_i\dot g_t+
 \langle\partial_{\rho}J_i,\dot{\hat\rho}_t\rangle+
 \partial_tJ_i.
\]
因此，即使第一项因负梯度控制而下降，总代价也可能被移动目标、旧地图或群体重构推高。更重要的是，$J_i$ 下降仅表示模型内指标改善，不等同于找到受困者、完成搬运或满足用户意图；任务成功还要由外部验收变量和证据链判断。原稿中的标量势可以保留为 $J_i$ 的可视化名称，但其来源、单位、归一化和更新时间必须记录。

邻域一致性项可取 $u_i^{coh}=k_v\sum_{j\in\mathcal N_i}a_{ij}(v_j-v_i)+k_f e_i^{form}$。第一项减少可通信邻居间的速度差，第二项纠正经任务绑定定义的队形误差；二者依赖动态图 $G_t$，并不来自真实磁场。只有在联合连通、时延有界和增益满足稳定条件时，才可推导渐近或实用一致性。探索项 $u_i^{exp}=\sigma_i\xi_i$ 可由零均值相关噪声或确定性覆盖序列产生，$\sigma_i$ 的单位必须与控制输入一致。这里的“温度” $T_i$ 定义为无量纲探索强度参数，经标定映射 $\sigma_i=f_T(T_i)$ 后才作用于执行器；它不是焦耳能耗、信息熵或设备热温度。耦合量 $\eta_i$ 同样是无量纲调度增益，用于改变一致性和任务权重，而不是自然常数。

离散实现采用步长 $\Delta t_k=t_{k+1}-t_k$，并在控制周期内零阶保持命令：
\[
 v_i^{k+1}=v_i^k+\frac{\Delta t_k}{m_i}
 \bigl(u_i^{safe,k}+u_i^{task,k}+u_i^{coh,k}+u_i^{exp,k}+d_i^k\bigr),
 \qquad r_i^{k+1}=r_i^k+\Delta t_k v_i^{k+1}.
\]
步长上限应由闭环带宽、最大相对速度和制动距离共同确定；网络包到达频率不能被默认等同于控制频率。消息过期时，控制器把相应边权衰减至零，并转入只依赖本地感知的保守模式。宏观层可通过调节 $T$ 与 $\eta$ 形成“分散探索、柔性协作、刚性编队”三个工作区间，但切换需要迟滞、驻留时间和安全投影，不能用两个旋钮承诺任意场景中的瞬时相变。可验证指标至少包括覆盖率、任务完成率、最小间距、连通分量数、单位任务能耗、消息新鲜度与控制饱和比例；这些量共同刻画性能，任何单一“场强”都不足以替代它们。

这一模型保留了原章关于斥力、目标牵引、队形一致性、探索扰动和三类工作模态的机制问题，同时撤回了不必要的物理等同。若任务地图光滑、通信较稳定且群体同质，梯度与一致性控制可给出有效近似；若环境存在离散门锁、单向通道或强异构能力，纯连续势函数会制造局部极小或错误平均，此时应由符号规划器、任务拍卖或CCM类协调模块提供离散约束，再由本节动力学执行。反例是所有机器人在门口形成对称拥堵：继续增大目标增益会提高推挤而非通过率，正确动作是改变任务绑定、给通行权并暂时降低局部耦合。动力学模型的价值正在于暴露这种可诊断失配，而不是用漂亮的轨迹替代任务证据。
